\subsection{Objetivos de diseño}

El diseño propuesto busca satisfacer cuatro objetivos simultáneos. En primer lugar, reducir la huella de memoria del vector de estado, que crece
exponencialmente con el número de qubits. En segundo lugar, preservar la exactitud numérica del simulador, evitando pérdidas de precisión en las
amplitudes complejas. En tercer lugar, impedir que cada compuerta obligue a reconstruir el vector de estado completo en memoria densa. Finalmente,
se requiere que la sobrecarga temporal introducida por la compresión permanezca controlada y pueda ser ajustada mediante parámetros de diseño.

Estos objetivos conducen a una separación clara entre la lógica del simulador y la estrategia de almacenamiento. La lógica cuántica continúa
operando sobre amplitudes complejas, compuertas y mediciones, mientras que la representación física del vector puede variar entre una forma densa
y una forma comprimida por bloques, siempre que ambas expongan el mismo comportamiento observable.
